\begin{abstract}

Commercial buildings can lower their electricity costs by committing ahead of time to demand response (DR) programs to shed their power draw during grid stress events, earning a capacity payment for committing to a firm cap on their power draw below their usual peak, and incurring a penalty for any overshoot. The DR programs provide incentive in addition to the time-of-use energy charges and monthly peak-aware demand charges. Most locations also provide Electric vehicle (EV) charging facilities. The EVs belonging to employees or part of the commercial fleet, parked at the workplace, are a natural resource to back the DR commitments, while also reducing the monthly peak, but they are not a reliable resource. The EVs arrive and depart on their users' schedules, consume uncertain amounts of energy, and their available charge on any day depends on their off-site usage and past charging decisions. The building usually commits a predetermined time ahead of an event, before the EVs and their energy requirements are known. We present P-CONSENT, a framework that turns this long-horizon problem, with uncertain and self-interested users into a committed demand-response resource through three coordinated decision policies: (i) a DR commitment decision policy that determines the optimal event peak using forecasts of next-day system conditions; (ii) a user negotiation policy, triggered upon user arrival, that incentivizes users, using the DR incentives and demand charge avoidance, to flexibly adjust their charging requirements and parking duration, and (iii) a charging policy that accounts for cross-day user behavior to achieve long-term demand charge reduction. Our policy prevents users from gaming the system, ensures that every payment is bounded by the value the user creates, and ensures voluntary participation. Our work is validated on operational data from a commercial building, and we demonstrate that the framework earns net demand-response revenue while lowering user charging costs and meeting its firm commitments.
\end{abstract}
